\worksheettitle
  {Ноль и разрядный состав}
  {\modulemark{M02} Учебный модуль \textbullet\ полный маршрут}

\studentnameblock

\begin{notebox}[Чему вы научитесь]
  Объяснять роль нуля в записи числа, раскладывать число на разрядные
  слагаемые и собирать число по такой сумме.
\end{notebox}

\begin{checkpointbox}[Вход: связь с M01]
  В числе \(81\,804\) две цифры \(8\).
  \begin{enumerate}[label=\textbf{\arabic*.}]
    \item Назовите их значения: \answerblank[24mm] и \answerblank[24mm].
    \item В каких разрядах стоят нули? \answerblank[65mm].
  \end{enumerate}
\end{checkpointbox}

\section{Зачем внутри числа нужен ноль}

Сравните записи \(40\,508\) и \(458\).

\begin{fillbox}[Пример с опорой: заполните разряды]
\begin{center}
\renewcommand{\arraystretch}{1.45}
\begin{tabular}{c|ccccc}
  \toprule
  разряд & дес. тысяч & ед. тысяч & сотни & десятки & единицы \\
  \midrule
  цифра в \(40\,508\) & 4 & \answerblank[9mm] & 5 & \answerblank[9mm] & 8 \\
  значение & 40 000 & \answerblank[13mm] & 500 & \answerblank[13mm] & 8 \\
  \bottomrule
\end{tabular}
\end{center}
\end{fillbox}

\newpage
\begin{practicebox}[Сравните значения цифр]
  \begin{enumerate}[label=\textbf{\arabic*.}]
    \item Значение цифры \(4\) в числе \(40\,508\): \answerblank[22mm],
      а в числе \(458\): \answerblank[22mm].
    \item Значение цифры \(5\) в числе \(40\,508\): \answerblank[22mm],
      а в числе \(458\): \answerblank[22mm].
    \item Почему числа \(40\,508\) и \(458\) различны?
      \writinglines{1}
  \end{enumerate}
\end{practicebox}

\begin{notebox}[Правило]
  Ноль показывает, что единиц данного разряда нет. Благодаря нулю можно
  однозначно определить разряды и значения остальных цифр.
\end{notebox}

\begin{practicebox}[Что произойдёт, если удалить ноль]
  Из числа \(60\,205\):
  \begin{enumerate}[label=\textbf{\arabic*.}]
    \item удалили ноль в разряде единиц тысяч. Получилось \answerblank[26mm];
    \item удалили ноль в разряде десятков. Получилось \answerblank[26mm].
  \end{enumerate}
  Сравните значения цифр до и после удаления нуля.
  \begin{enumerate}[label=\textbf{\arabic*.}]
    \item В обоих случаях значение цифры \(6\) изменилось с
      \answerblank[19mm] до \answerblank[19mm], то есть уменьшилось на
      \answerblank[22mm].
    \item Во втором случае значение цифры \(2\) изменилось с
      \answerblank[16mm] до \answerblank[16mm], то есть уменьшилось на
      \answerblank[18mm].
    \item Какое из вычисленных изменений больше? \answerblank[42mm].
  \end{enumerate}
\end{practicebox}

\section{Разложение на разрядные слагаемые}

\begin{notebox}[Разобранный пример и правило]
  Число можно представить в виде суммы разрядных слагаемых. Слагаемые,
  соответствующие нулевым цифрам, обычно не записывают:
  \[
    47\,305=40\,000+7\,000+300+5.
  \]
\end{notebox}

\begin{fillbox}[Заполните пропуски]
  \begin{enumerate}[label=\textbf{\arabic*.}]
    \item \(53\,204=50\,000+\answerblank[18mm]+200+\answerblank[16mm].\)
    \item \(80\,507=80\,000+\answerblank[18mm]+7.\)
    \item \(405\,060=400\,000+\answerblank[18mm]+60.\)
  \end{enumerate}
\end{fillbox}

\begin{practicebox}[Разложите самостоятельно]
  \begin{enumerate}[label=\textbf{\arabic*.}]
    \item \(63\,204=\answerblank[90mm].\)
    \item \(8\,090=\answerblank[90mm].\)
    \item \(50\,006=\answerblank[90mm].\)
    \item \(304\,080=\answerblank[90mm].\)
    \item \(600\,007=\answerblank[90mm].\)
    \item \(701\,040=\answerblank[90mm].\)
  \end{enumerate}
\end{practicebox}

\begin{reflectionbox}[Объясните]
  Почему в разложении числа \(304\,080\) нет слагаемых, соответствующих
  разрядам десятков тысяч, сотен и единиц?
  \writinglines{2}
\end{reflectionbox}

\begin{practicebox}[Разложите и соберите обратно]
  Разложите число \(70\,305\) на разрядные слагаемые. Затем сложите полученные
  слагаемые и убедитесь, что получилось исходное число.
  \writinglines{2}
\end{practicebox}

\section{Сборка числа по разрядным слагаемым}

\begin{fillbox}[Пример с опорой: заполните все разряды]
  Соберите число \(80\,000+700+5\).
  \begin{center}
  \renewcommand{\arraystretch}{1.45}
  \begin{tabular}{c|ccccc}
    \toprule
    разряд & дес. тысяч & ед. тысяч & сотни & десятки & единицы \\
    \midrule
    значение & 80 000 & \answerblank[11mm] & 700 & \answerblank[11mm] & 5 \\
    цифра & 8 & \answerblank[9mm] & 7 & \answerblank[9mm] & 5 \\
    \bottomrule
  \end{tabular}
  \end{center}
  Получилось число \answerblank[32mm].
\end{fillbox}

\begin{practicebox}[Запишите числа]
  \begin{enumerate}[label=\textbf{\arabic*.}]
    \item \(50\,000+3\,000+200+4=\answerblank[30mm].\)
    \item \(80\,000+700+5=\answerblank[30mm].\)
    \item \(300\,000+40\,000+8=\answerblank[30mm].\)
    \item \(600\,000+20+1=\answerblank[30mm].\)
    \item \(900\,000+4\,000+60=\answerblank[30mm].\)
    \item \(700\,000+5=\answerblank[30mm].\)
  \end{enumerate}
\end{practicebox}

\section{Разбираемся в ошибках}

\begin{warningbox}[Исправьте записи]
  Для каждой записи объясните ошибку и запишите верный результат.
  \begin{enumerate}[label=\textbf{\arabic*.}]
    \item \(80\,000+700+5=875.\)
    \item \(8\,090=8\,000+900.\)
    \item \(50\,204=50\,000+2\,000+200+4.\)
    \item \(300\,000+4\,000+8=34\,008.\)
  \end{enumerate}
  \writinglines{4}
\end{warningbox}

\section{Применяем способ}

\begin{practicebox}[Составьте числа по условиям]
  \begin{enumerate}[label=\textbf{\arabic*.}]
    \item Запишите пятизначное число, в котором нет единиц тысяч и десятков.
      Число: \answerblank[35mm].
    \item Запишите шестизначное число, разложение которого содержит ровно три
      ненулевых разрядных слагаемых. Число: \answerblank[35mm].
    \item Разложите оба составленных числа и проверьте выполнение условий.
      \writinglines{2}
  \end{enumerate}
\end{practicebox}

\begin{practicebox}[Вставьте нули]
  Вставьте в запись \(428\) один или два нуля так, чтобы получить:
  \begin{enumerate}[label=\textbf{\arabic*.}]
    \item четырёхзначное число: \answerblank[30mm];
    \item пятизначное число: \answerblank[30mm].
  \end{enumerate}
  Назовите значения цифр \(4\), \(2\) и \(8\) в каждом полученном числе.
  \writinglines{1}
\end{practicebox}

\begin{checkpointbox}[Контрольная точка]
  Выполните без подсказки.
  \begin{enumerate}[label=\textbf{\arabic*.}]
    \item \(90\,307=\answerblank[86mm].\)
    \item \(500\,040=\answerblank[86mm].\)
    \item \(30\,000+400+8=\answerblank[30mm].\)
    \item \(600\,000+2\,000+70=\answerblank[30mm].\)
    \item Почему в первом разложении нет слагаемых, соответствующих разрядам
      единиц тысяч и десятков?
      \writinglines{1}
  \end{enumerate}
\end{checkpointbox}

\begin{reflectionbox}[Проверяю себя]
  \choicebox\ Я объясняю, какие разряды обозначены нулями.

  \choicebox\ Я раскладываю число, не записывая нулевые слагаемые.

  \choicebox\ Я сохраняю все разряды при сборке числа по сумме.
\end{reflectionbox}

\begin{notebox}[Дальнейший маршрут]
  Если при сборке потерян разряд, цифре приписано неверное значение или роль
  нуля не объяснена, выполните M02-R. Если способ понятен, но нужна тренировка,
  выполните дополнительную практику или M02-H. Если контроль выполнен
  устойчиво, переходите к M03; M02\(\star\) можно выполнить дополнительно.
\end{notebox}
